\documentclass[11pt]{article}
\usepackage[T1]{fontenc}
\usepackage{lmodern}
\usepackage[margin=1in]{geometry}
\usepackage{amsmath,amssymb,amsthm,mathtools}
\usepackage{microtype}
\usepackage[hidelinks]{hyperref}
\hypersetup{pdftitle={An Exact Semidefinite Lift of a Nonspectrahedral Hyperbolicity Cone},pdfauthor={OpenAI}}
\newtheorem{theorem}{Theorem}[section]
\newtheorem{lemma}[theorem]{Lemma}
\newtheorem{proposition}[theorem]{Proposition}
\newtheorem{corollary}[theorem]{Corollary}
\theoremstyle{remark}
\newtheorem{remark}[theorem]{Remark}
\DeclareMathOperator{\tr}{tr}
\DeclareMathOperator{\adj}{adj}
\DeclareMathOperator{\Rea}{Re}
\DeclareMathOperator{\range}{range}
\newcommand{\R}{\mathbb R}
\newcommand{\C}{\mathbb C}
\newcommand{\Sym}{\mathbb S}
\title{An Exact Semidefinite Lift of a\protect\\Nonspectrahedral Hyperbolicity Cone}
\author{OpenAI}
\date{October 5, 2026}
\begin{document}
\maketitle
\begin{abstract}
We construct an exact semidefinite lift of the explicit nonspectrahedral
hyperbolicity cone in twenty-three variables defined in the companion
paper. The lift is a homogeneous real symmetric pencil of size $100$
with $307$ auxiliary variables and represents the entire closed cone,
including every point with singular $X$. Thus, although this cone has no
semidefinite representation in its original coordinates, it admits one
when auxiliary variables are allowed. The stated sizes are not claimed
to be minimal.
\end{abstract}
\section{Introduction}
\label{sec:introduction}

Let $\Sym^n$ denote the real symmetric $n\times n$ matrices. A
\emph{spectrahedron} is a set defined by a finite affine linear matrix
inequality; a \emph{spectrahedral shadow} is a linear projection of such
a set. Thus auxiliary variables may make a semidefinite representation
possible even when no representation exists in the original variables.
We study this distinction for one explicit hyperbolicity cone.

For a homogeneous real polynomial $p$ with $p(e)>0$, hyperbolicity with
respect to $e$ means that $t\mapsto p(te-x)$ has only real roots for
every $x$. We use the closed-cone convention
\[
 K(p,e)=\{x:\text{every root of }t\mapsto p(te-x)
                    \text{ is nonnegative}\}.
\]
In particular, zero roots are allowed.

The convexity of hyperbolicity cones goes back to
G\aa rding~\cite{garding1959}. In three variables, the
Helton--Vinnikov theorem yields definite symmetric determinantal
representations, and hence spectrahedral
cones~\cite{helton-vinnikov2007,lewis-parrilo-ramana2005}.
Allowing auxiliary variables leads to a broader representation question.
Netzer and Sanyal proved semidefinite liftability when every nonzero
boundary point is a smooth point of the defining hyperbolic
polynomial~\cite{netzer-sanyal2015}. More recently, Scheiderer obtained
second-order cone lifts under a Nash-smooth boundary
hypothesis~\cite{scheiderer2025}. Our result instead concerns an explicit
cone and gives its lift by a bounded-degree algebraic construction.

\subsection{The cone and the main result}
For $y=(y_1,y_2,y_3)\in\R^3$, define
\begin{equation}\label{eq:Q}
 Q(y)=\begin{pmatrix}
 y_1^2+y_2^2&-y_1y_2&-y_1y_3\\
 -y_1y_2&y_2^2+y_3^2&-y_2y_3\\
 -y_1y_3&-y_2y_3&y_3^2+y_1^2
 \end{pmatrix}.
\end{equation}
The associated biquadratic form $z^TQ(y)z$ is the coefficient-one
Choi--Lam form~\cite[Section~4]{choi-lam1977}.
For $a\in\R$, $r\in\R^3$, and $T\in\Sym^3$, set
\begin{equation}\label{eq:Phi}
 \Phi_y\!\begin{pmatrix}a&r^T\\r&T\end{pmatrix}
 =\begin{pmatrix}
 \tr(Q(y)T)&-r^TQ(y)\\
 -Q(y)r&aQ(y)
 \end{pmatrix}.
\end{equation}
This is a linear map from $\Sym^4$ to itself for each fixed $y$.
On the space $\mathcal V=\Sym^4\times\Sym^4\times\R^3$, define
\begin{equation}\label{eq:p}
 p(X,Z,y)=\det\!\left((\det X)Z-\Phi_y(\adj X)\right),
 \qquad e=(I_4,I_4,0),
\end{equation}
where $\adj X$ denotes the classical adjugate. Each matrix entry inside
this determinant is homogeneous of degree five, and $p(e)=1$; hence
$p$ is homogeneous of degree twenty. We write $K=K(p,e)$.
The slice representation in Section~\ref{sec:slices} will in particular
verify that $p$ is hyperbolic with respect to $e$.

Identify $\mathcal V$ with $\R^{23}$ by listing the upper triangular
entries of $X$, then those of $Z$, each in lexicographic order of
$(i,j)$, followed by $y_1,y_2,y_3$.

\begin{theorem}\label{thm:main}
There exist fixed matrices $A_1,\ldots,A_{23},B_1,\ldots,B_{307}
\in\Sym^{100}$ such that, for every $x=(X,Z,y)\in\mathcal V$,
\[
 x\in K\quad\Longleftrightarrow\quad
 \text{there exists }v\in\R^{307}\text{ with }
 \sum_{\nu=1}^{23}x_\nu A_\nu+
 \sum_{j=1}^{307}v_jB_j\succeq0.
\]
The equivalence holds on the entire closed cone, including points
with singular $X$.
\end{theorem}

The sizes in Theorem~\ref{thm:main} are explicit bounds, with no
minimality assertion. The matrices are specified by a finite
coefficient rule in Section~\ref{sec:lift}.

The cone in \eqref{eq:p} was constructed in the companion
manuscript~\cite{parent}, whose Theorem~1.1 proves that it has no
homogeneous semidefinite representation in its original twenty-three
coordinates. This also excludes affine representations. Indeed, if a
cone containing $0$ equals $\{x:A_0+L(x)\succeq0\}$, where $L$ is
linear, then $A_0\succeq0$. For each point $x$ of the cone and $t>0$,
we have $A_0/t+L(x)\succeq0$; letting $t\to\infty$ gives $L(x)\succeq0$.
Conversely, $L(x)\succeq0$ implies $A_0+L(x)\succeq0$; thus the same
cone has the homogeneous representation $L(x)\succeq0$.
We use the companion result only for this comparison: the proof of
Theorem~\ref{thm:main} below is self-contained. The companion also
establishes hyperbolicity and, in Proposition~3.2, describes the part of
$K$ with $X\succ0$. Our slice calculation recovers those two facts as
part of the lift construction.

\subsection{Proof strategy}
The starting point is a family of planes in $\R^3$:
\[
 y=Y_\theta(\eta,\xi)
   =\eta(\cos\theta,\sin\theta,0)+\xi(0,0,1).
\]
They cover all values of $y$. On each plane, the matrix $Q(y)$ has a
Gram factorization linear in $(\eta,\xi)$ and trigonometric of degree
three in $\theta$. This gives a $20\times20$ linear pencil $L_\theta$
whose positive semidefinite locus is exactly the corresponding slice
of $K$, including its boundary.

The angular parameter does not itself furnish a finite linear matrix
inequality. We instead introduce a real linear functional on a
finite-dimensional space of trigonometric polynomials with coefficients
linear in the slice variables. We require it to be nonnegative on
$q^TL_\theta q$ for every vector $q$ of trigonometric polynomials of
degree at most two. These requirements form one finite matrix
inequality. Evaluation at a single angle provides a feasible functional
for every point of $K$.

This use of a finite matrix to encode positivity against a space of
test polynomials belongs to the moment and sum-of-squares approach to
semidefinite representations~\cite{lasserre2001,lasserre2009}.
Here exactness must hold at one fixed degree.

The main issue is the converse: an arbitrary feasible functional need
not be evaluation at an angle, or integration against a measure.
The degree-two matching identity in Section~\ref{sec:matching} supplies
enough test vectors to recover the matrix inequalities defining each
slice. Its construction rotates a fixed complex Gram factor by a
$2\times2$ unitary matrix. The rotation preserves the Gram matrix and
reduces a cubic trigonometric pairing to a first harmonic. Matching
its value and first derivative then identifies the pairing exactly.
This bounded-degree construction is the principal additional algebraic
ingredient of the lift. Section~\ref{sec:lift} applies it to prove
exactness for every feasible functional and completes the proof of
Theorem~\ref{thm:main}.

\section{Exact semidefinite representations on angular slices}
\label{sec:slices}

We first construct a matrix pencil on each plane $y=Y_\theta(\eta,\xi)$.
Its determinant will agree with $p$ on that plane. This will identify
cone membership with positive semidefiniteness even when $X$ is singular.

For $\theta\in\mathbb R$, set
\[
 c=\cos\theta,\qquad s=\sin\theta,\qquad w=e^{i\theta},
 \qquad d(\theta)=(c,s,0),\qquad k=(0,0,1),
\]
and write
\[
 Y_\theta(\eta,\xi)=\eta d(\theta)+\xi k
 \qquad (\eta,\xi\in\mathbb R).
\]
These planes cover $\mathbb R^3$.  We first factor $Q$ on each plane,
using complex notation to keep the formulas short.  For a complex matrix
$F=(f_1,f_2)\in\mathbb C^{3\times2}$, define its realification by
\[
 \widetilde F=(\operatorname{Re}f_1,\operatorname{Im}f_1,
                 \operatorname{Re}f_2,\operatorname{Im}f_2).
\]
Thus $\widetilde F\widetilde F^T=\operatorname{Re}(FF^*)$, where $*$
denotes conjugate transpose.  Define
\begin{equation}\label{eq:U}
 U(\theta)=\frac1{\sqrt2}
 \begin{pmatrix}
 -w^{-1}&-w\\
 sw^{-2}&sw^2\\
 cw^{-2}&-cw^2
 \end{pmatrix},\qquad
 V_0=\frac1{\sqrt2}
 \begin{pmatrix}0&0\\ i&i\\1&1\end{pmatrix},
 \qquad F_\theta(\eta,\xi)=\eta U(\theta)+\xi V_0.
\end{equation}

\begin{lemma}[Angular Gram factorization]\label{lem:gram}
For every $\theta,\eta,\xi\in\mathbb R$,
\begin{equation}\label{eq:gram}
 \widetilde F_\theta(\eta,\xi)
 \widetilde F_\theta(\eta,\xi)^T
 =Q\bigl(Y_\theta(\eta,\xi)\bigr).
\end{equation}
In particular, $Q(y)\succeq0$ for every $y\in\mathbb R^3$.
\end{lemma}

\begin{proof}
The three coefficients in the quadratic expansion of
$\operatorname{Re}(F_\theta F_\theta^*)$ are
\begin{align*}
 \operatorname{Re}(UU^*)&=
 \begin{pmatrix}1&-sc&0\\-sc&s^2&0\\0&0&c^2\end{pmatrix},
 &\operatorname{Re}(V_0V_0^*)&=
 \begin{pmatrix}0&0&0\\0&1&0\\0&0&1\end{pmatrix},\\
 \operatorname{Re}(UV_0^*+V_0U^*)&=
 \begin{pmatrix}0&0&-c\\0&0&-s\\-c&-s&0\end{pmatrix}.
\end{align*}
For the $(2,3)$ entry of the last matrix, the multiplication gives
$s\cos(2\theta)-c\sin(2\theta)=-s$.
The sum of these matrices with coefficients $\eta^2,\xi^2,\eta\xi$
is exactly $Q(\eta c,\eta s,\xi)$.  Every real $y$ has this form,
so the factorization also proves positive semidefiniteness.
\end{proof}


\subsection{The slice pencil}
For a real column $r\in\R^3$ and a real vector
$u=(u_0,u')\in\R\times\R^3$, define
\begin{equation}\label{eq:BJ}
 B(r)=\begin{pmatrix}0&r^T\\-r&0_{3\times3}\end{pmatrix},
 \qquad J_u=\begin{pmatrix}u'&-u_0I_3\end{pmatrix}.
\end{equation}
Then $B(r)u=J_u^Tr$. If $r_1,\ldots,r_4$ are the columns of
$\widetilde F_\theta(\eta,\xi)$, direct multiplication and
Lemma~\ref{lem:gram} give
\begin{equation}\label{eq:Phi-factor}
 \sum_{j=1}^4 B(r_j)^TWB(r_j)
       =\Phi_{Y_\theta(\eta,\xi)}(W)
       \qquad(W\in\Sym^4).
\end{equation}
Indeed, for $W=\left(\begin{smallmatrix}a&v^T\\v&T\end{smallmatrix}\right)$,
the individual summands have blocks
$r_j^TTr_j$, $-r_j^Tv\,r_j^T$, $-r_jr_j^Tv$, and $a r_jr_j^T$.
Summing uses $\sum_jr_jr_j^T=Q(Y_\theta(\eta,\xi))$.
Consequently, for every $y\in\R^3$,
\begin{equation}\label{eq:Phi-quadratic}
 u^T\Phi_y(W)u=\tr(J_uWJ_u^TQ(y)).
\end{equation}
Here we used that the angular planes cover $\R^3$.

Write $\mathcal X,\mathcal Z\in\Sym^4$ for the matrix variables on a
slice, and define
\begin{equation}\label{eq:slice-pencil}
 L_\theta(\mathcal X,\mathcal Z,\eta,\xi)
 =\begin{pmatrix}
 I_4\otimes\mathcal X&\mathcal B_\theta(\eta,\xi)\\
 \mathcal B_\theta(\eta,\xi)^T&\mathcal Z
 \end{pmatrix},
 \qquad
 \mathcal B_\theta(\eta,\xi)=
 \begin{pmatrix}B(r_1)\\B(r_2)\\B(r_3)\\B(r_4)\end{pmatrix}.
\end{equation}
For fixed $\theta$, this is a homogeneous real symmetric linear pencil
of size twenty. Its coefficients are trigonometric polynomials in
$\theta$ of degree at most three.

\begin{lemma}\label{lem:slice}
For every $\theta,\eta,\xi\in\R$ and
$\mathcal X,\mathcal Z\in\Sym^4$,
\begin{equation}\label{eq:slice-equivalence}
 (\mathcal X,\mathcal Z,Y_\theta(\eta,\xi))\in K
 \quad\Longleftrightarrow\quad
 L_\theta(\mathcal X,\mathcal Z,\eta,\xi)\succeq0.
\end{equation}
In particular, $p$ is hyperbolic with respect to $e$. If
$\mathcal X\succ0$, the conditions in \eqref{eq:slice-equivalence}
are also equivalent to
$\mathcal Z\succeq\Phi_{Y_\theta(\eta,\xi)}(\mathcal X^{-1})$.
\end{lemma}

\begin{proof}
For invertible $\mathcal X$, the block determinant formula and
\eqref{eq:Phi-factor} show that
\begin{align*}
 \det L_\theta
 &= (\det\mathcal X)^4
    \det\!\left(\mathcal Z-
        \Phi_{Y_\theta(\eta,\xi)}(\mathcal X^{-1})\right)\\
 &=p(\mathcal X,\mathcal Z,Y_\theta(\eta,\xi)).
\end{align*}
The second equality uses
$\adj\mathcal X=(\det\mathcal X)\mathcal X^{-1}$ and linearity of
$\Phi_y$ in its matrix argument. Both sides are polynomials in the
slice variables, so the identity also holds for singular $\mathcal X$.
Moreover, $L_\theta(I_4,I_4,0,0)=I_{20}$. Thus
\begin{equation}\label{eq:roots}
 p\!\left(te-(\mathcal X,\mathcal Z,Y_\theta(\eta,\xi))\right)
 =\det\!\left(tI_{20}-L_\theta(\mathcal X,\mathcal Z,\eta,\xi)\right).
\end{equation}
The roots are the real eigenvalues of this symmetric pencil evaluated
at the specified point. Since every $y$ lies on some angular plane,
this proves hyperbolicity. Nonnegativity of all these eigenvalues,
including zero eigenvalues, proves \eqref{eq:slice-equivalence}.
For $\mathcal X\succ0$, the final assertion is the Schur complement
criterion applied to \eqref{eq:slice-pencil}.
\end{proof}

\section{A degree-two matching identity}\label{sec:matching}

Write
\[
 Q(y,z)=\tfrac12\bigl(Q(y+z)-Q(y)-Q(z)\bigr)
\]
for the symmetric bilinear polarization of $Q$, so that $Q(y,y)=Q(y)$.

For $S\in\mathbb S^3$ and $F,H\in\mathbb C^{3\times2}$, put
\[
 \langle F,H\rangle_S
 =\operatorname{Re}\operatorname{tr}(F^*SH)
 =\operatorname{tr}(\widetilde F^T S\widetilde H).
\]
This is a symmetric real bilinear form.  A complex trigonometric
polynomial of degree at most $j$ is a linear combination of
$e^{i\ell\theta}$ for $-j\leq\ell\leq j$.

The finite lift will test $L_\theta$ against trigonometric polynomial
vectors built from the columns of a matrix $\widetilde C(\theta)$.
A constant Gram matrix for $\widetilde C$ makes the coefficient of
$\mathcal X$ in the quadratic contribution of the top blocks independent
of $\theta$. Requiring
$\langle F_\theta,C\rangle_S=\operatorname{tr}(SQ(Y_\theta,z))$
makes the mixed contribution linear in $Y_\theta$.
The next lemma achieves both properties with entries of
degree at most two. The matrix $S$ in the pairing may be indefinite.

\begin{lemma}[Degree-two matching]\label{lem:matching}
Let $S\in\mathbb S^3$ and let $z\in\mathbb R^3$ satisfy
$(z_1,z_2)\ne(0,0)$.  There is a matrix
$C(\theta)\in\mathbb C^{3\times2}$ whose entries are trigonometric
polynomials of degree at most two such that, for all
$\theta,\eta,\xi\in\mathbb R$,
\begin{equation}\label{eq:matching}
 \begin{aligned}
 \widetilde C(\theta)\widetilde C(\theta)^T&=Q(z),\\
 \langle F_\theta(\eta,\xi),C(\theta)\rangle_S
 &=\operatorname{tr}\bigl(SQ(Y_\theta(\eta,\xi),z)\bigr).
 \end{aligned}
\end{equation}
\end{lemma}

\begin{proof}
Write $z=\tau d(\alpha)+\rho k$, where $\tau>0$, and set
\[
 w_0=e^{i\alpha},\qquad R_0=\tau U(\alpha)+\rho V_0,
 \qquad N=SR_0.
\]
By Lemma~\ref{lem:gram}, $\widetilde R_0\widetilde R_0^T=Q(z)$.
We will set $C(\theta)=R_0M(\theta)$ with $M(\theta)$ unitary, so
that this Gram matrix remains fixed.  We choose the rotation to make
$\langle V_0,C\rangle_S$ constant and
$\langle U(\theta),C(\theta)\rangle_S$ a linear combination of
$\cos\theta$ and $\sin\theta$.  After those frequency bounds are
established, values and a derivative at $\alpha$ will identify the
pairings.

\smallskip
\noindent\emph{Choosing the unitary rotation.}
Write $N_j$ for the $j$th row of $N$ and introduce the two-component row
\[
 v=N_3-iN_2.
\]
Subscripts $+$ and $-$ denote the first and second components of a row.
Set
\[
 g=v_++\overline{v_-},\qquad
 h=v_+-\overline{v_-},\qquad
 \kappa(\theta)=(w_0/w)^2,\qquad
 \delta=|g|^2+|h|^2.
\]
We seek a matrix of the form
\[
 M(\theta)=\begin{pmatrix}a&b\\-\overline b&\overline a\end{pmatrix}.
\]
Writing $v^M=vM$, we impose the two conditions
\begin{equation}\label{eq:rotation-pairings}
 v^M_++\overline{v^M_-}=g,
 \qquad
 v^M_+-\overline{v^M_-}=h\kappa.
\end{equation}
The sum on the left controls the pairing with $V_0$, while the difference
multiplies $w^3$ in the pairing with $U$.  These conditions therefore
keep the first pairing constant and lower the cubic term of the second
to frequency one, since $w^3\kappa=w_0^2w$.  The expansion below will
also bound its remaining frequencies.

The two left sides of \eqref{eq:rotation-pairings} are
$ga+\overline h\,\overline b$ and
$ha-\overline g\,\overline b$, respectively.  Thus, for $\delta>0$,
we choose $a(\theta)$ and $b(\theta)$ by solving
\begin{equation}\label{eq:rotation-system}
 \begin{pmatrix}g&\overline h\\h&-\overline g\end{pmatrix}
 \begin{pmatrix}a\\\overline b\end{pmatrix}
 =\begin{pmatrix}g\\h\kappa\end{pmatrix}.
\end{equation}
The coefficient matrix $A$ satisfies $A^*A=\delta I_2$, so the
system has a unique solution.  Explicitly,
\begin{equation}\label{eq:rotation-coefficients}
 a=\frac{|g|^2+|h|^2\kappa}{\delta},\qquad
 \overline b=\frac{gh(1-\kappa)}{\delta}.
\end{equation}
Since the right side of \eqref{eq:rotation-system} has squared norm
$\delta$, we have $|a|^2+|b|^2=1$.
For $\delta=0$, take $a=1$ and $b=0$; then $g=h=0$ and the same
system is satisfied.  In both cases $M$ is unitary,
$M(\alpha)=I_2$, and $a,\overline b$ involve only the
frequencies $0$ and $-2$.  Thus $C=R_0M$ has degree at most two and
$C(\alpha)=R_0$.  Moreover,
\[
 \widetilde C\widetilde C^T
 =\operatorname{Re}(R_0MM^*R_0^*)
 =\operatorname{Re}(R_0R_0^*)=Q(z),
\]
which proves the first identity in \eqref{eq:matching}.

\smallskip
\noindent\emph{Reducing the frequencies of the pairings.}
The formula for $V_0$ and \eqref{eq:rotation-pairings} yield
\[
 \sqrt2\,\langle V_0,C\rangle_S
 =\operatorname{Re}(v^M_++v^M_-)=\operatorname{Re}g,
\]
so this pairing is constant.

Set $m=N_3+iN_2$, $n=N_1$, and write $m^M=mM$, $n^M=nM$.
Expanding $c$ and $s$ in powers of $w$, the formula for $U$ gives
\begin{align*}
 \sqrt2\,\langle U(\theta),C(\theta)\rangle_S
 =\operatorname{Re}\bigl[{}
 &\tfrac12w^3v^M_+
 +w(\tfrac12m^M_+-n^M_+)\\
 &-\tfrac12w^{-3}v^M_-
 +w^{-1}(-\tfrac12m^M_--n^M_-)
 \bigr].
\end{align*}
Inside the real part, we may replace each negative-power term by its
complex conjugate.  The resulting expression is
\begin{equation}\label{eq:frequency-reduction}
 \operatorname{Re}\left[
 \frac{w^3}{2}(v^M_+-\overline{v^M_-})
 +w\left(\frac{m^M_+}{2}-n^M_+
          -\frac{\overline{m^M_-}}{2}-\overline{n^M_-}\right)
 \right].
\end{equation}
The first term has frequency one because
$w^3\kappa=w_0^2w$.  In the second term, the expression in parentheses
has only frequencies $0$ and $-2$: for any fixed row $r$,
\[
 (rM)_+=r_+a-r_-\overline b,
 \qquad
 \overline{(rM)_-}=\overline{r_+}\,\overline b+\overline{r_-}a.
\]
Consequently \eqref{eq:frequency-reduction} is a real linear
combination of $\cos\theta$ and $\sin\theta$, as required.

\smallskip
\noindent\emph{Identifying the pairings.}
At $\theta=\alpha$, polarization of Lemma~\ref{lem:gram} within the
plane spanned by $d(\alpha)$ and $k$ gives
\begin{align}
 \langle V_0,R_0\rangle_S
 &=\operatorname{tr}\bigl(SQ(k,z)\bigr),\label{eq:matching-value-v}\\
 \langle U(\alpha),R_0\rangle_S
 &=\operatorname{tr}\bigl(SQ(d(\alpha),z)\bigr).
 \label{eq:matching-value-u}
\end{align}
Both sides of the desired identity
$\langle V_0,C(\theta)\rangle_S=\operatorname{tr}(SQ(k,z))$
are constant, so \eqref{eq:matching-value-v} proves it for every
$\theta$.

For the pairing with $U$, both
$\langle U(\theta),C(\theta)\rangle_S$ and
$\operatorname{tr}(SQ(d(\theta),z))$ are linear combinations of
$\cos\theta$ and $\sin\theta$.  They agree at $\alpha$ by
\eqref{eq:matching-value-u}; it remains to compare their derivatives
there.  The fixed Gram matrix of $C$ gives
\[
 \langle C(\theta),C(\theta)\rangle_S
 =\operatorname{tr}(SQ(z)),
 \qquad
 \langle R_0,C'(\alpha)\rangle_S=0.
\]
Constancy of the pairing with $V_0$ also gives
$\langle V_0,C'(\alpha)\rangle_S=0$.  Since
$R_0=\tau U(\alpha)+\rho V_0$ and $\tau>0$, it follows that
\begin{equation}\label{eq:matching-moving-factor}
 \langle U(\alpha),C'(\alpha)\rangle_S=0.
\end{equation}
Next differentiate the identity
\[
 \langle\tau U(\theta)+\rho V_0,
          \tau U(\theta)+\rho V_0\rangle_S
 =\operatorname{tr}\bigl(SQ(\tau d(\theta)+\rho k)\bigr),
\]
which follows from Lemma~\ref{lem:gram}.  Evaluating at $\alpha$ and
dividing by $2\tau$ gives
\[
 \langle U'(\alpha),R_0\rangle_S
 =\operatorname{tr}\bigl(SQ(d'(\alpha),z)\bigr).
\]
Together with \eqref{eq:matching-moving-factor}, this is the required
derivative identity.  A real linear combination of $\cos\theta$ and
$\sin\theta$ is determined by its value and derivative at a single
angle.  Hence the pairing with $U$ has the desired value for all
$\theta$.  Linearity in $\eta,\xi$ completes the proof.
\end{proof}

\section{A finite semidefinite lift}\label{sec:lift}

We now replace the angular family of pencils by one finite matrix condition.
Evaluation on an angular slice will provide a feasible lift for each point of
$K$.  Conversely, the functions supplied by Lemma~\ref{lem:matching} will
turn positivity of the finite matrix into the inequalities that characterize
the slice, including its singular boundary.

\subsection{The matrix condition}

Let $\mathcal T_j$ be the real vector space of trigonometric polynomials of
degree at most $j$, with basis
\[
  1,\cos\theta,\sin\theta,\ldots,\cos(j\theta),\sin(j\theta).
\]
Let $\mathcal E$ be the space of real functions of
$(\theta,\mathcal X,\mathcal Z,\eta,\xi)$ that are linear homogeneous in
the $22$ coordinates of
$(\mathcal X,\mathcal Z,\eta,\xi)\in
\mathbb S^4\times\mathbb S^4\times\mathbb R^2$, with coefficients in
$\mathcal T_7$.  Thus $\dim\mathcal E=22\cdot15=330$.
Our lifting variable is a real linear functional
$\Lambda:\mathcal E\to\mathbb R$; on vectors and matrices we apply it
entrywise.  For an output point $(X,Z,y)$, impose
\begin{equation}\label{eq:output}
  \Lambda(\mathcal X)=X,\qquad
  \Lambda(\mathcal Z)=Z,\qquad
  \Lambda\bigl(Y_\theta(\eta,\xi)\bigr)=y.
\end{equation}
We also require
\begin{equation}\label{eq:moment}
  \Lambda\bigl(q(\theta)^T
  L_\theta(\mathcal X,\mathcal Z,\eta,\xi)q(\theta)\bigr)\geq0
  \qquad\text{for every }q\in(\mathcal T_2)^{20}.
\end{equation}
The expression to which $\Lambda$ is applied lies in $\mathcal E$:
the entries of $L_\theta$ have trigonometric degree at most three, so its
product with two entries of $q$ has degree at most seven.

Condition~\eqref{eq:moment} is a single finite semidefinite condition,
using the finite test-space construction familiar from moment
relaxations~\cite{lasserre2001,lasserre2009}.
Indeed, set $(f_1,\ldots,f_5)=
(1,\cos\theta,\sin\theta,\cos2\theta,\sin2\theta)$ and define
$H(\Lambda)\in\mathbb S^{100}$ by
\begin{equation}\label{eq:moment-matrix}
  H(\Lambda)_{(i,a),(j,b)}
    =\Lambda\bigl(f_a f_b (L_\theta)_{ij}\bigr),
  \qquad 1\leq i,j\leq20,\quad 1\leq a,b\leq5.
\end{equation}
If $q_i=\sum_a t_{i,a}f_a$, the left side of~\eqref{eq:moment} is
$t^TH(\Lambda)t$.  Hence~\eqref{eq:moment} is equivalent to
$H(\Lambda)\succeq0$.  Every entry of this real symmetric matrix is linear
homogeneous in the $330$ coordinates of $\Lambda$.

\begin{samepage}
\begin{proposition}\label{prop:lift}
For $(X,Z,y)\in\mathbb S^4\times\mathbb S^4\times\mathbb R^3$, the
following are equivalent:
\begin{enumerate}
  \item $(X,Z,y)\in K$;
  \item there is a real linear functional $\Lambda:\mathcal E\to\mathbb R$
  satisfying~\eqref{eq:output} and $H(\Lambda)\succeq0$.
\end{enumerate}
\end{proposition}
\end{samepage}

\begin{proof}
Suppose first that $(X,Z,y)\in K$.  Choose
$\theta_0,\eta_0,\xi_0$ with
$y=Y_{\theta_0}(\eta_0,\xi_0)$, and let $\Lambda$ be evaluation at
$(\theta_0,X,Z,\eta_0,\xi_0)$.  The output equations hold, and
Lemma~\ref{lem:slice} gives
$L_{\theta_0}(X,Z,\eta_0,\xi_0)\succeq0$, which implies
\eqref{eq:moment}.  This evaluation functional is a finite feasible
witness even when $X$ is singular.

For the converse, suppose that $\Lambda$ satisfies the stated conditions.
Taking constant $q$ supported on one of the four top blocks of $L_\theta$
in~\eqref{eq:moment} gives $X\succeq0$.
We next obtain a family of scalar inequalities from which the remaining
slice conditions will follow.

Fix $u\in\mathbb R^4$, $D\in\mathbb S^4$, and
$z\in\mathbb R^3$ with $(z_1,z_2)\neq(0,0)$, and apply
Lemma~\ref{lem:matching} with
$S=J_uDJ_u^T\in\mathbb S^3$.  Neither $D$ nor $S$ is required to be
positive semidefinite.  Let $b_1(\theta),\ldots,b_4(\theta)$ be
the columns of the resulting real matrix $\widetilde C(\theta)$.
Choose the five four-dimensional blocks of $q$ to be
\[
  q(\theta)=
  \begin{pmatrix}
    -DJ_u^Tb_1(\theta)\\
    \vdots\\
    -DJ_u^Tb_4(\theta)\\
    u
  \end{pmatrix}.
\]
The degree bound in Lemma~\ref{lem:matching} puts this vector in
$(\mathcal T_2)^{20}$.  Write $r_j$ for the columns of
$\widetilde F_\theta(\eta,\xi)$, as in the definition of $L_\theta$.
Using $B(r_j)u=J_u^Tr_j$, the top blocks and the cross terms of
$q^TL_\theta q$ are respectively
\begin{align*}
  \sum_{j=1}^4 b_j^T J_uD\mathcal X DJ_u^Tb_j
    &=\operatorname{tr}\bigl(J_uD\mathcal X DJ_u^TQ(z)\bigr),\\
  -2\sum_{j=1}^4 b_j^TJ_uDJ_u^Tr_j
    &=-2\langle F_\theta(\eta,\xi),C(\theta)\rangle_S\\
    &=-2\operatorname{tr}\bigl(J_uDJ_u^T
        Q(Y_\theta(\eta,\xi),z)\bigr).
\end{align*}
These identities hold before applying $\Lambda$.  The first identity
removes the angular dependence from the coefficient of $\mathcal X$;
the last expression is linear in $Y_\theta(\eta,\xi)$.
Consequently, applying~\eqref{eq:output} and~\eqref{eq:moment} gives
\begin{equation}\label{eq:certificate}
  u^TZu+
  \operatorname{tr}\bigl(J_uDXDJ_u^TQ(z)\bigr)
  -2\operatorname{tr}\bigl(J_uDJ_u^TQ(y,z)\bigr)\geq0.
\end{equation}
For fixed $u,D$, the left side is a polynomial in $z$, so the inequality
extends by continuity to every $z\in\mathbb R^3$.  Only this scalar
inequality is extended; no continuity of the functions $C$ is needed.
Thus~\eqref{eq:certificate} holds for all $u\in\mathbb R^4$,
$D\in\mathbb S^4$, and $z\in\mathbb R^3$.

If $X$ were positive definite, setting $z=y$ and $D=X^{-1}$ in
\eqref{eq:certificate} would give $Z\succeq\Phi_y(X^{-1})$.
For singular $X$, the same family of tests also supplies the necessary
range condition.  Fix slice parameters
$y=Y_{\theta_0}(\eta_0,\xi_0)$ and, for the rest of the proof, let
$r_j$ denote the columns of
$\widetilde F_{\theta_0}(\eta_0,\xi_0)$ and set $B_j=B(r_j)$.
For $v\in\ker X$, take $D=t vv^T$ and $z=y$ in
\eqref{eq:certificate}, where $t\in\mathbb R$ is arbitrary.  Since
$DXD=0$, this gives
\[
  u^TZu-2t\,(J_uv)^TQ(y)(J_uv)\geq0
  \qquad(t\in\mathbb R).
\]
Its coefficient of $t$ must vanish.  The Gram factorization of $Q(y)$
therefore yields
\[
  0=(J_uv)^TQ(y)(J_uv)
    =\sum_{j=1}^4(v^TB_ju)^2.
\]
As $u$ and $v\in\ker X$ are arbitrary, it follows that
\begin{equation}\label{eq:range}
  \operatorname{range}B_j\subseteq(\ker X)^\perp
    =\operatorname{range}X\qquad(1\leq j\leq4).
\end{equation}

Let $X^\dagger$ be the symmetric matrix that inverts the positive
eigenvalues of $X$ and is zero on $\ker X$.  In particular,
$X^\dagger XX^\dagger=X^\dagger$.
Taking $D=X^\dagger$ and $z=y$ in~\eqref{eq:certificate}, and using
\eqref{eq:Phi-quadratic} and~\eqref{eq:Phi-factor}, gives
\begin{equation}\label{eq:singular-schur}
  Z\succeq\Phi_y(X^\dagger)
    =\sum_{j=1}^4B_j^TX^\dagger B_j.
\end{equation}
We now verify positivity of the slice pencil directly by completing
a square using~\eqref{eq:range} and~\eqref{eq:singular-schur}.  For
$a_1,\ldots,a_4,u\in\mathbb R^4$, its quadratic form equals
\begin{align*}
  &\sum_{j=1}^4a_j^TXa_j
    +2\sum_{j=1}^4a_j^TB_ju+u^TZu\\
  &\quad=\sum_{j=1}^4
    (a_j+X^\dagger B_ju)^TX(a_j+X^\dagger B_ju)
    +u^T\left(Z-\sum_{j=1}^4B_j^TX^\dagger B_j\right)u
    \geq0.
\end{align*}
Here the cross terms agree because $XX^\dagger B_j=B_j$ by
\eqref{eq:range}.  Thus
$L_{\theta_0}(X,Z,\eta_0,\xi_0)\succeq0$, and
Lemma~\ref{lem:slice} gives $(X,Z,y)\in K$.
\end{proof}

The converse has used an arbitrary feasible functional $\Lambda$.
In particular, the proof requires no representation of $\Lambda$ by a
measure and no closure operation on the projected feasible set.

\subsection{The size of the lift}

\begin{proof}[Proof of Theorem~\ref{thm:main}]
Choose the displayed trigonometric basis of $\mathcal T_7$ and the
upper-triangular coordinate functions of $\mathcal X$ and $\mathcal Z$,
together with $\eta,\xi$, to specify the $330$ coordinates of $\Lambda$.
The output equations~\eqref{eq:output} prescribe precisely the following
$23$ distinct coordinates:
\begin{gather*}
  \Lambda(\mathcal X_{ij}),\quad
  \Lambda(\mathcal Z_{ij})\qquad(1\leq i\leq j\leq4),\\
  \Lambda(\eta\cos\theta),\quad
  \Lambda(\eta\sin\theta),\quad
  \Lambda(\xi).
\end{gather*}
They are independent members of the chosen coordinate system on
$\mathcal E^*$.  Substitute the corresponding entries of $X,Z,y$ for
these coordinates in~\eqref{eq:moment-matrix}, and use the remaining
$330-23=307$ coordinates as auxiliary variables.  The result is a fixed
real symmetric $100\times100$ matrix pencil, homogeneous in the output
and auxiliary variables.  Proposition~\ref{prop:lift} proves that its
projection is exactly $K$.  Hence the requested representation has
$N=100$, $m=307$, and zero constant term.
\end{proof}

\bibliographystyle{alpha}
\bibliography{references}
\end{document}
